\documentclass[10pt,a4paper]{amsart}
\usepackage[margin=19mm]{geometry}
\usepackage{amsmath,amssymb,mathtools}
\usepackage[scr=rsfs,cal=euler]{mathalfa}
\usepackage{xfrac}
\usepackage[unicode,hidelinks,bookmarks=false]{hyperref}
\newcommand{\ie}{i.e.\ }
\newcommand{\cf}{cf.\ }
\newcommand{\resp}{respectively\ }
\newcommand{\Subsection}{\subsection}
\providecommand{\nobreakdash}{\nobreak\hbox{-}}
\setlength{\emergencystretch}{2em}
\multlinegap0pt
\usepackage{bm}
\usepackage{bbm}
\usepackage{dsfont}
\usepackage{xspace}
\usepackage{cancel}
\usepackage{mathtools}
\usepackage{amsthm}
\makeatletter
\@ifundefined{thm}{\newtheorem{thm}{Theorem}}{}
\makeatother
\title[Derivative Type Mapping Theorem for the Interpolative Berind]{Derivative Type Mapping Theorem for the Interpolative Berinde Weak Contraction in Metric Spaces with Application}
\author{Clement Boateng Ampadu}
\date{Source: arXiv:2604.14157v1 (CC BY 4.0)}
\begin{document}
\maketitle

\noindent\emph{Source and licence.} Derivative Type Mapping Theorem for the Interpolative Berinde Weak Contraction in Metric Spaces with Application. Version arXiv:2604.14157v1 (CC BY 4.0), distributed under the Creative Commons Attribution 4.0 International licence, as recorded in the arXiv licence metadata for this exact version and shown on the abstract page https://arxiv.org/abs/2604.14157v1. Full rights notice: licenses/arxiv-2604.14157-CC-BY-4.0.txt.\par\smallskip

\noindent\textit{Editorial scope.} Counted unit: the Thm whose source label is \texttt{None} in the article (source0.tex lines 136--137, source label \texttt{None}), delivered together with the article's own complete original proof of it (source0.tex lines 139--168, one \texttt{proof} environment of 30 lines). No mathematics is written, repaired or re-derived: statement and proof are the article's own text line for line. Required context retained uncounted: none. Every definition, display and label of those lines is retained unchanged. Omitted from this package and not redistributed: 1--135, 138--138, 169--235. Nothing inside a retained excerpt is omitted or altered: every retained body is delivered exactly as the article's lines are written, so no omission receipt of any kind accompanies this package. Citations: the retained excerpts carry no citation command at all, so no bibliography entry of the article is transcribed and no reference is redistributed. Reference adaptation: none was needed and none was made; every reference inside the delivered unit resolves inside it, so no unresolved reference or citation is produced. Printed slips kept literal: no printed word, symbol, hypothesis or number of the counted statement or of its proof was altered. Layout and class: the article's own class is \texttt{article} with packages including fontenc, inputenc, lmodern, authblk, amssymb, amsfonts, xy, enumerate, mathrsfs, amsthm, amsmath, tikz, mathtools, float; the wrapper is the lane's pinned \texttt{amsart} editorial wrapper at 10pt on A4, which restates the theorem-like environments the retained text needs and adds the article's own macro definitions that the retained text needs (none), with extra wrapper packages (bm, bbm, dsfont, xspace, cancel, mathtools). The delivered numbering is the unit's own. Assets: no retained span contains a figure environment, a graphics-inclusion command, a PDF-page inclusion, a table or a TikZ picture; no image file is copied into this package, the package declares no asset dependency and no image file is redistributed with it. Licence: version arXiv:2604.14157v1 of this article is distributed under the Creative Commons Attribution 4.0 International licence, as recorded in the arXiv licence metadata for this exact version and shown on the abstract page https://arxiv.org/abs/2604.14157v1. A full rights notice is delivered at licenses/arxiv-2604.14157-CC-BY-4.0.txt. Characters: every retained excerpt is pure ASCII, so the delivered engine, XeLaTeX with its default Latin Modern text font, reports no missing character.
\begin{thm} Let $(X,d)$ be a metric space, and $T:X\mapsto X$ be an interpolative Berinde weak operator of the derivative type. Suppose that $X$ is $d$-bounded. Let $\varphi: \mathbb{R}^{+}\mapsto \mathbb{R}^{+}$ be a function such that $\frac{d\varphi}{dt}\bigg|_{t=\epsilon} >0$ for each $\epsilon>0$. Then $T$ has a unique fixed point $x^{*}\in X$, such that for each $x\in X$, $\lim_{n\to\infty} T^{n}x=x^{*}$
\end{thm}

\begin{proof} Let $x_0\in X$ be chosen arbitrarily. Define the sequence $\{x_n\}_{n=0}^{\infty}$ recursively by $x_{n+1}=Tx_n$ for all $n$. Letting $x=x_{n-1}$ and $y=x_n$ in the contractive definition of the theorem we have
\begin{align*}
\frac{d\varphi}{dt}\bigg|_{t=d(x_n,x_{n+1})} & \leq \lambda \bigg[\frac{d\varphi}{dt}\bigg|_{t=d(x_{n-1},x_n)}\bigg]^{\alpha} \bigg[\frac{d\varphi}{dt}\bigg|_{t=d(x_{n-1},x_n)}\bigg]^{1-\alpha} \\
&= \lambda \frac{d\varphi}{dt}\bigg|_{t=d(x_{n-1},x_n)}
\end{align*}
By induction we obtain
\begin{align*}
\frac{d\varphi}{dt}\bigg|_{t=d(x_n,x_{n+1})} & \leq \lambda^{n} \frac{d\varphi}{dt}\bigg|_{t=d(x_0,x_1)} \\
& \leq \lambda^{n} \frac{d\varphi}{dt}\bigg|_{t=\delta_d(X)}
\end{align*}
where $d(x_0,x_1)\leq \delta_d(X)$, and $\delta_d(X)$ is stated as in Definition 1.3. Thus making use of the above inequality and the triangle inequality we have for all $n\geq 0$ and $m\geq 1$ that 
\begin{align*}
\frac{d\varphi}{dt}\bigg|_{t=d(x_n,x_{n+m})}&\leq \frac{d\varphi}{dt}\bigg|_{t=d(x_n,x_{n+1})}+\frac{d\varphi}{dt}\bigg|_{t=d(x_{n+1},x_{n+2})}+\cdots+ \frac{d\varphi}{dt}\bigg|_{t=d(x_{n+m-1},x_{n+m})} \\
&\leq \frac{d\varphi}{dt}\bigg|_{t=\delta_d(X)} (\lambda^{n}+\lambda^{n+1}+\cdots+\lambda^{n+m-1})\\
&=\frac{d\varphi}{dt}\bigg|_{t=\delta_d(X)} \lambda ^{n} \frac{1-\lambda^{m}}{1-\lambda}\\
&\leq \frac{d\varphi}{dt}\bigg|_{t=\delta_d(X)} \frac{\lambda^{n}}{1-\lambda}~\rightarrow 0~as~n,m \rightarrow \infty
\end{align*}
Since $\frac{d\varphi}{dt}\bigg|_{t=d(x_n,x_{n+m})} \rightarrow 0$, by the condition on $\varphi$, we have $d(x_n,x_{n+m})\rightarrow 0$. This shows that
$\{x_n\}$ is a Cauchy sequence in the complete metric space $(X,d)$, ensuring its convergence to a point $x^{*}\in X$. To confirm that $x^{*}$ is a fixed point, we
substitute $x=x^{*}$ and $y=x_n$ in the contractive definition of the theorem, then we have
$$ \frac{d\varphi}{dt}\bigg|_{t=d(Tx^{*},x_{n+1})} \leq \lambda \bigg[\frac{d\varphi}{dt}\bigg|_{t=d(x^{*},x_n)}\bigg]^{\alpha} \bigg[\frac{d\varphi}{dt}\bigg|_{t=d(x^{*},Tx^{*})}\bigg]^{1-\alpha}$$ Letting $n\rightarrow \infty$ in the above inequality, we have
$$\frac{d\varphi}{dt}\bigg|_{t=d(Tx^{*},x^{*})}\leq \lambda \bigg[\frac{d\varphi}{dt}\bigg|_{t=0}\bigg]^{\alpha} \bigg[\frac{d\varphi}{dt}\bigg|_{t=d(x^{*},Tx^{*})}\bigg]^{1-\alpha}$$ The condition on $\varphi$ gives $\frac{d\varphi}{dt}\bigg|_{t=0}=0$. So from the above inequality we have
$$\frac{d\varphi}{dt}\bigg|_{t=d(Tx^{*},x^{*})} \leq 0$$ which is a contradiction. Therefore by the condition on $\varphi$ we obtain
$$\frac{d\varphi}{dt}\bigg|_{t=d(Tx^{*},x^{*})}=0$$ this leads to $d(Tx^{*},x^{*})=0$ or $Tx^{*}=x^{*}$. We now prove that $T$ has a unique fixed point. Suppose this is not true, then there exists $x=Tx$, $y=Ty$, $x\neq y$, $d(x,y)>0$. Therefore using the contractive definition of the theorem, we have
\begin{align*}
\frac{d\varphi}{dt}\bigg|_{t=d(x,y)} &= \frac{d\varphi}{dt}\bigg|_{t=d(Tx,Ty)}\\
&\leq \lambda \bigg[\frac{d\varphi}{dt}\bigg|_{t=d(x,y)}\bigg]^{\alpha} \bigg[\frac{d\varphi}{dt}\bigg|_{t=0}\bigg]^{1-\alpha}
\end{align*}
The condition on $\varphi$ gives, $\frac{d\varphi}{dt}\bigg|_{t=0}=0$. Thus from the above inequality we have $$\frac{d\varphi}{dt}\bigg|_{t=d(x,y)}\leq 0$$ a contradiction. Again the condition on $\varphi$ gives, $\frac{d\varphi}{dt}\bigg|_{t=d(x,y)}=0$, which leads to $d(x,y)=0$ or $x=y$. This completes the proof.
\end{proof}

\end{document}
